\DocumentMetadata{
  pdfversion=2.0,
  pdfstandard=ua-2,
  lang=en-US,
  tagging=on,
  tagging-setup={math/setup=mathml-SE,tabsorder=structure}
}
\documentclass[11pt,letterpaper]{article}
\usepackage[margin=0.68in,includefoot]{geometry}
\usepackage{newcomputermodern}
\usepackage{amsmath,amscd,mathtools}
\providecommand{\square}{\mdlgwhtsquare}
\usepackage[hidelinks]{hyperref}
\hypersetup{
  pdftitle={Numerical Analysis PhD Exam, August 2019, Part 1},
  pdfauthor={University of Florida Department of Mathematics},
  pdfsubject={Graduate mathematics examination},
  pdfdisplaydoctitle=true,
  bookmarksopen=true
}
\setcounter{secnumdepth}{0}
\setlength{\parindent}{0pt}
\setlength{\parskip}{0.28em}
\setlength{\emergencystretch}{3em}
\setlength{\leftmargini}{2.15em}
\setlength{\leftmarginii}{2.65em}
\setlength{\leftmarginiii}{3em}
\renewcommand{\labelenumi}{\textbf{\theenumi.}}
\pagestyle{empty}
\raggedbottom
\allowdisplaybreaks
\newenvironment{examproblems}[1]{%
  \begin{enumerate}%
  \setcounter{enumi}{#1}%
  \setlength{\topsep}{0.35em}%
  \setlength{\partopsep}{0pt}%
  \setlength{\itemsep}{0.48em}%
  \setlength{\parsep}{0.18em}%
}{\end{enumerate}}
\newenvironment{examsubparts}{%
  \begin{enumerate}%
  \setlength{\topsep}{0.18em}%
  \setlength{\partopsep}{0pt}%
  \setlength{\itemsep}{0.16em}%
  \setlength{\parsep}{0.1em}%
}{\end{enumerate}}
\newenvironment{examinstructions}{%
  \par\begingroup\itshape
}{%
  \par\endgroup\vspace{0.18em}
}
\makeatletter
\renewcommand\section{\@startsection{section}{1}{0pt}%
  {-1.25ex plus -.25ex minus -.1ex}%
  {0.55ex plus .1ex}%
  {\normalfont\large\bfseries}}
\renewcommand{\@maketitle}{%
  \newpage\null\vskip -1.2em%
  {\footnotesize\bfseries
    \href{https://www.ufl.edu/}{University of Florida}, \href{https://math.ufl.edu/}{Department of Mathematics}\par}%
  \vskip 0.18em%
  \tagstructbegin{tag=Title}%
    {\LARGE\bfseries\href{https://gma.math.ufl.edu/exams/numerical-analysis-phd/}{Numerical Analysis PhD Exam}, August 2019, Part 1\par}%
  \tagstructend%
  \vskip 0.35em%
  \tagmcbegin{artifact}%
    \hrule height 1pt%
  \tagmcend%
  \vskip 0.55em%
}
\makeatother
\title{Numerical Analysis PhD Exam, August 2019, Part 1}
\author{}
\date{}
\begin{document}
\maketitle
\thispagestyle{empty}
\begin{examinstructions}
Do 4 (four) problems.
\end{examinstructions}
\begin{examproblems}{0}
\item
Let \(A \in \mathbb{C}^{m\times n}\).
\begin{examsubparts}
\item[\textnormal{(a)}]
Prove or give a counterexample: \(\|A\|_2 \leq \sqrt{\|A\|_\infty\|A\|_1}\). If you prove this, make sure to justify each nontrivial step.
\item[\textnormal{(b)}]
Prove or give a counterexample: \(\|A\|_2 \leq \|A\|_F\), where \(\|A\|_F\) is the Frobenius norm of~\(A\). If you prove this, make sure to justify each nontrivial step.
\end{examsubparts}
\item
This problem has two parts.
\begin{examsubparts}
\item[\textnormal{(a)}]
Prove that the inverse of an upper-triangular matrix is upper-triangular.
\item[\textnormal{(b)}]
Let \(A \in \mathbb{C}^{m\times n}\) with \(m>n\). Consider the least-squares problem of finding \(x \in \mathbb{C}^n\) that minimizes \(\|Ax-b\|\) in the 2-norm. Describe a method for solving the problem efficiently, and explain (and justify) why the normal equations should not be solved.
\end{examsubparts}
\item
Let \(A \in \mathbb{C}^{m\times n}\), with \(m\geq n\) and \(\operatorname{rank}(A)=p=n\geq 3\). Let \(a_1,a_2,\ldots\) denote the columns of~\(A\).
\begin{examsubparts}
\item[\textnormal{(a)}]
Using the classical Gram–Schmidt process, write out expressions for \(q_1,q_2,q_3\), the first three columns of~\(Q\) in the \(QR\) decomposition of~\(A\).
\item[\textnormal{(b)}]
Show the vector \(q_3\) found in part (a) is orthogonal to both \(q_1\) and \(q_2\).
\item[\textnormal{(c)}]
Write an expression for the first Householder reflector \(H_1\), used to find the \(QR\) decomposition of~\(A\). Show \(H_1\) is both unitary and Hermitian.
\end{examsubparts}
\item
Let \(A \in \mathbb{C}^{m\times m}\) be Hermitian.
\begin{examsubparts}
\item[\textnormal{(a)}]
Show that all eigenvalues of~\(A\) are real.
\item[\textnormal{(b)}]
Define the stationary iterative method (a.k.a. fixed point method) 
\[
x^{(k+1)}=Ax^{(k)}+b. \tag{1}
\]
 Suppose (1) has fixed-point~\(x\), namely~\(x\) satisfies \(x=Ax+b\). Show the iteration (1) converges to~\(x\) from any starting guess \(x^{(0)}\), that is \(x^{(k)}\to x\) as \(k\to\infty\), if and only if the eigenvalues \(\lambda_i\) of~\(A\) satisfy \(|\lambda_i|<1\), \(i=1,\ldots,m\). You may use the fact that Hermitian matrix~\(A\) is unitarily diagonalizable.
\end{examsubparts}
\item
Consider the matrix~\(A\) given by 
\[
\begin{pmatrix}
1 & -1 & 2 & 0\\
-1 & 4 & -1 & 1\\
2 & -1 & 6 & -2\\
0 & 1 & -2 & 4
\end{pmatrix}.
\]
 Suppose the eigenvalues of~\(A\) are all distinct (they are) and satisfy \(\lambda_1>\lambda_2>\lambda_3>\lambda_4\).
\begin{examsubparts}
\item[\textnormal{(a)}]
Show that~\(A\) is positive definite.
\item[\textnormal{(b)}]
Describe an algorithm that could be used to approximate \(\lambda_4\).
\item[\textnormal{(c)}]
Describe algorithms that could be used to approximate \(\lambda_2,\lambda_3\), and their eigenvectors.
\end{examsubparts}
\end{examproblems}
\end{document}
